\documentclass[reqno]{amsart}
\usepackage{mathptmx}      % use Times fonts if available on your TeX system
\usepackage{latexsym}
\usepackage{amssymb}
\usepackage{bm}
\usepackage[usenames,dvipsnames]{xcolor}
\usepackage{graphicx,pifont}
\usepackage{multicol,multirow}
\usepackage{ulem,cancel}
\usepackage{tikz}
\usetikzlibrary[patterns]
\usepackage{comment}
\usepackage{hyperref}
\usepackage{array}
\usepackage{caption}
\usepackage{subcaption}
\usepackage{txfonts}
\usepackage[shortlabels]{enumitem}
\usepackage{cleveref}
\crefname{equation}{}{}
\crefname{enumi}{}{}
\crefname{figure}{Figure}{Figure}
\crefname{subsection}{Subsection}{Subsections}
\crefname{lemma}{Lemma}{Lemma}
\crefname{theorem}{Theorem}{Theorem}
\crefname{proposition}{Proposition}{Proposition}
\crefname{section}{Section}{Section}

\newtheorem{theorem}{Theorem}[section]
\newtheorem{lemma}[theorem]{Lemma}
\newtheorem{proposition}{Proposition}[section]
\theoremstyle{definition}
\newtheorem{definition}[theorem]{Definition}
\newtheorem{example}[theorem]{Example}
\theoremstyle{remark}
\newtheorem{remark}[theorem]{Remark}
\numberwithin{equation}{section}
\numberwithin{figure}{section}
\newtheorem{conjecture}[theorem]{Conjecture}

\newcommand{\crossref}{{\color{red}ref here}}

\newcommand{\bs}[1]{\boldsymbol{#1}}
\newcommand{\bsb}[1]{{\mbox{\sffamily #1}}}
\newcommand{\bb}[1]{{\mbox{\sffamily \bfseries #1}}}
\newcommand{\oname}[1]{\textrm{#1}}
\newcommand{\abs}[1]{\left|#1\right|}
\newcommand{\nrm}[1]{\left|\left|#1\right|\right|}
\newcommand{\eqdef}{\stackrel{\mathrm{def}}{=\joinrel=}}
\newcommand{\ab}[1]{\langle#1\rangle}
\newcommand{\po}[1]{\frac{\partial}{\partial{#1}}}
\newcommand{\poh}[1]{\partial/\partial{#1}}
\newcommand{\pp}[2]{\frac{\partial {#1}}{\partial {#2}}}
\newcommand{\pph}[2]{\partial {#1}/\partial {#2}}
\newcommand{\termc}[1]{{\mbox{\sffamily {#1}}}}
\newcommand{\bt}[1]{{\boldsymbol{\mbox{\sffamily \bfseries {#1}}}}}
\newcommand{\phf}[1]{{#1}+1/2}
\newcommand{\mhf}[1]{{#1}-1/2}
\newcommand{\powth}{{\textrm{\scriptsize th}}}
\newcommand{\powrd}{{\textrm{\scriptsize rd}}}
\newcommand{\powst}{{\textrm{\scriptsize st}}}
\newcommand{\pownd}{{\textrm{\scriptsize nd}}}

\setlength{\unitlength}{1cm}
\setlength{\arraycolsep}{3pt} 
\setlength{\tabcolsep}{1pt}

\begin{document}

\title[Numerical Oncolytic Dynamics]{Finite Volume Modeling of Oncolytic Viral Therapy}
\author[X.~Zeng]{Xianyi Zeng}
\address{Department of Mathematics, Lehigh University, Bethlehem PA 18015, United States. \\    
         Tel.: +1-610-758-3745}
\email[X.~Zeng]{xyzeng@lehigh.edu}
\date{\today}

%\subjclass[2010]{65M12 \and 35L65}

%\keywords{
%}

%\begin{abstract}
%\end{abstract}

\maketitle

\section{Governing Equations of the Oncolytic Viral Therapy}
\label{sec:eqn}
{\bf The governing equations}.
The PDE for normalized cell number densities in radial coordinate is given by:
\begin{subequations}\label{eq:eqn_pde_cell}
  \begin{align}
    \label{eq:eqn_pde_cell_b}
    &\pp{B}{t} + \frac{1}{r^2}\po{r}\left[r^2BE\right] = (\lambda - \beta V - k_2\theta M_2)B\;, \\
    \label{eq:eqn_pde_cell_h}
    &\pp{H}{t} + \frac{1}{r^2}\po{r}\left[r^2HE\right] = \beta V B - (k_1M_1\theta+k_3M_2\theta+\delta)H\;, \\
    \label{eq:eqn_pde_cell_q}
    &\pp{Q}{t} + \frac{1}{r^2}\po{r}\left[r^2QE\right] = g(Q)Q-k_4M_3\theta Q\;, \\
    \label{eq:eqn_pde_cell_n}
    &\pp{N}{t} + \frac{1}{r^2}\po{r}\left[r^2NE\right] = \delta H + k_1M_1\theta H + k_2M_2\theta B + k_3M_2\theta H + k_4M_3\theta Q - \rho N\;, \\
    \label{eq:eqn_pde_cell_m1}
    &\pp{M_1}{t} + \frac{1}{r^2}\po{r}\left[r^2M_1\left(E+U_1\right)\right] = \left[s_1H\theta-c_1-f_1(H)\right]M_1 {\color{red} + h_1(r,t)}\;, \\
    \label{eq:eqn_pde_cell_m2}
    &\pp{M_2}{t} + \frac{1}{r^2}\po{r}\left[r^2M_2\left(E+U_2\right)\right] = \left[s_2(B+H)\theta-c_2-f_2(B,H)\right]M_2 {\color{red} + h_2(r,t)}\;, \\
    \label{eq:eqn_pde_cell_m3}
    &\pp{M_3}{t} + \frac{1}{r^2}\po{r}\left[r^2M_3\left(E+U_3\right)\right] = \left[s_3Q\theta-c_3-f_3(Q)\right]M_3 {\color{red} + h_3(r,t)}\;, \\
    \label{eq:eqn_pde_cell_m1e}
    &\pp{M_1^0}{t} + \frac{1}{r^2}\po{r}\left[r^2M_1^0E\right] = f_1(H)M_1-c_1^0M_1^0\;, \\
    \label{eq:eqn_pde_cell_m2e}
    &\pp{M_2^0}{t} + \frac{1}{r^2}\po{r}\left[r^2M_2^0E\right] = f_2(B,H)M_2-c_2^0M_2^0\;, \\
    \label{eq:eqn_pde_cell_m3e}
    &\pp{M_3^0}{t} + \frac{1}{r^2}\po{r}\left[r^2M_3^0E\right] = f_3(Q)M_3-c_3^0M_3^0\;, \\
    \label{eq:eqn_pde_cell_s}
    &{\color{red}\pp{S}{t} + \frac{1}{r^2}\po{r}\left[r^2SE\right] = -\kappa(S-S_0) - h_1(r,t) - h_2(r,t) - h_3(r,t)}\;,
  \end{align}
\end{subequations}
where $\theta$ is the (constant) total cell number density; the velocities $U_1$, $U_2$, $U_3$, $E$ are:
\begin{subequations}\label{eq:eqn_pde_vel}
  \begin{align}
    \label{eq:eqn_pde_vel_u}
    &U_1 = \alpha_1\theta\pp{H}{r}\;,\quad U_2 = \hat{\alpha}_2\theta\pp{B}{r}+\alpha_2\theta\pp{H}{r}\;,\quad U_3 = \alpha_3\theta\pp{Q}{r}\;, \\
    \notag
    &\frac{1}{r^2}\po{r}\left\{r^2\left[E+M_1U_1+M_2U_2+M_3U_3\right]\right\} \\
    \label{eq:eqn_pde_vel_e}
    &\quad = \lambda B + g(Q)Q-\rho N +s_1H\theta M_1+s_2(B+H)\theta M_2+s_3Q\theta M_3 \\
    \notag
    &\qquad -c_1M_1-c_2M_2-c_3M_3-c_1^0M_1-c_2^0M_2-c_3^0M_3{\color{red}-\kappa(S-S_0)}\;;
  \end{align}
\end{subequations}
and lastly the virus number density $V$ is determined by the diffusion equation:
\begin{equation}\label{eq:eqn_pde_virus}
  \pp{V}{t} = \frac{1}{r^2}\po{r}\left(D_1r^2\pp{V}{r}\right) + b\delta\theta H-\gamma V\;.
\end{equation}
Note that~\cref{eq:eqn_pde_vel_e} is obtained by adding all equations of~\cref{eq:eqn_pde_cell} and substitute $B+H+Q+N+M_1+M_2+M_3+M_1^0+M_2^0+M_3^0{\color{red}+S}=1$ in the left hand side.
Meanwhile we do not specify the function $g(Q)$ as it will hold that $Q\equiv0$ throughout the computation; and the other three functions are given by:
\begin{equation}\label{eq:eqn_pde_efun}
  f_1(H) = \frac{d_1p_H H}{K_1/\theta+p_HH}\,,\ 
  f_2(B,H) = \frac{d_2(p_BB+p_HH)}{K_2/\theta+p_BB+p_HH}\,,\ 
  f_3(Q) = \frac{d_3p_QQ}{K_3/\theta+p_QQ}\,.
\end{equation}
{\color{red}
The three functions $h_1(r,t)$, $h_2(r,t)$, and $h_3(r,t)$ represent the supply of immune cells from blood veins; for now, they are set to be:
\begin{equation}\label{eq:eqn_pde_supply}
  %h_1(r,t) = h_2(r,t) = h_3(r,t) = \delta_0(t-T_M)\sum_{l=1}^L\frac{n_M}{t_M\theta}\delta_0(r-r_l)\;,
  h_1(r,t) = h_2(r,t) = h_3(r,t) = \delta_0(t-T_M)\sum_{l=1}^L\frac{n_M}{\theta}\delta_0(r-r_l)\;,
\end{equation}
where $r_l=5l\textrm{mm}$, $l=1,2,\cdots$ are prescribed locations inside the tumor, $n_M$ is the number of vein spots per $\textrm{mm}^2$ on a spherical surface (one immune cell pops out from one such spot at one time, we assume), 
%$t_M$ is the time for one immune cell to pop out of the vein, 
and $T_M$ is the time at which the immune cells are injected into the vein.
(Note that due to radial symmetry, this function indicates the entire surface with radius $r_l$ secretes immune cells uniformly.)
}
Lastly, all equations are defined on the domain $r\in[0,\;R(t)]$, where $R(t)$ is the radius of the tumor that is determined by:
\begin{equation}\label{eq:eqn_pde_rad}
  \dot{R}(t) = E(R(t),t)\;.
\end{equation}

{\bf Parameters}. 
In the list below, we list all parameters in the model as well as their corresponding variables in the program.
\begin{itemize}
  \item {\color{cyan}$\theta=10^6\;\textrm{cell}\cdot\textrm{mm}^{-3}$}; accessed by {\color{green}\texttt{tdCoef->theta}}.
  \item {\color{cyan}$\lambda=0.042\;\textrm{h}^{-1}$}; accessed by {\color{green}\texttt{["B"]["grow\_pos"]}}.
  \item {\color{cyan}$\beta=0.1\times10^{-6}\;\textrm{mm}^3\cdot\textrm{virus}^{-1}\cdot\textrm{h}^{-1}$}; accessed by {\color{green}\texttt{["B"]["coef\_vi\_neg"]}}.
  \item {\color{cyan}$k_2=1.2\times10^{-6}\;\textrm{mm}^3\cdot\textrm{cell}^{-1}\cdot\textrm{h}^{-1}$}; accessed by {\color{green}\texttt{["B"]["coef\_m2\_neg"]}}.
  \item {\color{cyan}$k_1=1.2\times10^{-6}\;\textrm{mm}^3\cdot\textrm{cell}^{-1}\cdot\textrm{h}^{-1}$}; accessed by {\color{green}\texttt{["H"]["coef\_m1\_neg"]}}.
  \item {\color{cyan}$k_3=1.2\times10^{-6}\;\textrm{mm}^3\cdot\textrm{cell}^{-1}\cdot\textrm{h}^{-1}$}; accessed by {\color{green}\texttt{["H"]["coef\_m2\_neg"]}}.
  \item {\color{cyan}$\delta=0.042\;\textrm{h}^{-1}$}; accessed by {\color{green}\texttt{["H"]["grow\_neg"]}}.
  \item {\color{cyan}$k_4=1.2\times10^{-6}\;\textrm{mm}^3\cdot\textrm{cell}^{-1}\cdot\textrm{h}^{-1}$}; accessed by {\color{green}\texttt{["Q"]["coef\_m3\_neg"]}}.
  \item {\color{cyan}$\rho=0.67\;\textrm{h}^{-1}$}; accessed by {\color{green}\texttt{["N"]["grow\_neg"]}}.
  \item {\color{cyan}$s_1=0.5\times10^{-6}\;\textrm{mm}^3\cdot\textrm{cell}^{-1}\cdot\textrm{h}^{-1}$}; accessed by {\color{green}\texttt{["M1"]["coef\_h\_pos"]}}.
  \item {\color{cyan}$c_1=0.0075\;\textrm{h}^{-1}$}; accessed by {\color{green}\texttt{["M1"]["grow\_neg"]}}.
  \item {\color{cyan}$p_H=7\times10^6$}; accessed by {\color{green}\texttt{["M1"]["mm\_slope"]}} and {\color{green}\texttt{["M2"]["mm\_slope\_h"]}}.
  \item {\color{cyan}$K_1=10\times10^{15}$}; accessed by {\color{green}\texttt{["M1"]["mm\_half"]}}.
  \item {\color{cyan}$\alpha_1=1.0\times10^{-6}\;\textrm{mm}^5\cdot\textrm{cell}^{-1}\cdot\textrm{h}^{-1}$}; accessed by {\color{green}\texttt{["M1"]["vel\_grad\_h"]}}.
  \item {\color{cyan}$d_1=0.0083\;\textrm{h}^{-1}$}; accessed by {\color{green}\texttt{["M1"]["grow\_exhaust\_neg"]}}.
  \item {\color{cyan}$s_2=0.5\times10^{-6}\;\textrm{mm}^3\cdot\textrm{cell}^{-1}\cdot\textrm{h}^{-1}$}; accessed by {\color{green}\texttt{["M2"]["coef\_bh\_pos"]}}.
  \item {\color{cyan}$c_2=0.0075\;\textrm{h}^{-1}$}; accessed by {\color{green}\texttt{["M2"]["grow\_neg"]}}.
  \item {\color{cyan}$p_B=7\times10^6$}; accessed by {\color{green}\texttt{["M2"]["mm\_slope\_b"]}}.
  \item {\color{cyan}$K_2=10\times10^{15}$}; accessed by {\color{green}\texttt{["M2"]["mm\_half"]}}.
  \item {\color{cyan}$\hat{\alpha}_2=1.0\times10^{-6}\;\textrm{mm}^5\cdot\textrm{cell}^{-1}\cdot\textrm{h}^{-1}$}; accessed by {\color{green}\texttt{["M2"]["vel\_grad\_b"]}}.
  \item {\color{cyan}$\alpha_2=1.0\times10^{-6}\;\textrm{mm}^5\cdot\textrm{cell}^{-1}\cdot\textrm{h}^{-1}$}; accessed by {\color{green}\texttt{["M2"]["vel\_grad\_h"]}}.
  \item {\color{cyan}$d_2=0.0083\;\textrm{h}^{-1}$}; accessed by {\color{green}\texttt{["M2"]["grow\_exhaust\_neg"]}}.
  \item {\color{cyan}$s_3=0.5\times10^{-6}\;\textrm{mm}^3\cdot\textrm{cell}^{-1}\cdot\textrm{h}^{-1}$}; accessed by {\color{green}\texttt{["M3"]["coef\_q\_pos"]}}.
  \item {\color{cyan}$c_3=0.0075\;\textrm{h}^{-1}$}; accessed by {\color{green}\texttt{["M3"]["grow\_neg"]}}.
  \item {\color{cyan}$p_Q=7\times10^6$}; accessed by {\color{green}\texttt{["M3"]["mm\_slope"]}}.
  \item {\color{cyan}$K_3=10\times10^{15}$}; accessed by {\color{green}\texttt{["M3"]["mm\_half"]}}.
  \item {\color{cyan}$\alpha_3=1.0\times10^{-6}\;\textrm{mm}^5\cdot\textrm{cell}^{-1}\cdot\textrm{h}^{-1}$}; accessed by {\color{green}\texttt{["M3"]["vel\_grad\_q"]}}.
  \item {\color{cyan}$d_3=0.0083\;\textrm{h}^{-1}$}; accessed by {\color{green}\texttt{["M3"]["grow\_exhaust\_neg"]}}.
  \item {\color{cyan}$c_1^0=0.0075\;\textrm{h}^{-1}$}; accessed by {\color{green}\texttt{["M1E"]["grow\_neg"]}}.
  \item {\color{cyan}$c_2^0=0.0075\;\textrm{h}^{-1}$}; accessed by {\color{green}\texttt{["M2E"]["grow\_neg"]}}.
  \item {\color{cyan}$c_3^0=0.0075\;\textrm{h}^{-1}$}; accessed by {\color{green}\texttt{["M3E"]["grow\_neg"]}}.
  \item {\color{red} $S_0=0.15$}; accessed by {\color{green}\texttt{["S"]["equilibrium"]}}.
  \item {\color{red} $\kappa=10^2\;\textrm{h}^{-1}$}; accessed by {\color{green}\texttt{["S"]["elasticity"]}}.
  \item {\color{red} $n_M=1\;\textrm{cell}\cdot\textrm{mm}^{-2}$}; accessed by {\color{green}\texttt{["M1"]["popup\_num"]}}, {\color{green}\texttt{["M2"]["popup\_num"]}}, and {\color{green}\texttt{["M3"]["popup\_num"]}}.
  %\item {\color{red} $t_M=10^{-3}\;\textrm{h}$}; accessed by {\color{green}\texttt{["M1"]["popup\_time"]}}, {\color{green}\texttt{["M2"]["popup\_time"]}}, and {\color{green}\texttt{["M3"]["popup\_time"]}}.
  \item {\color{red} $T_M=36\;\textrm{h}$}; accessed by {\color{green}\texttt{["M1"]["inject\_time"]}}, {\color{green}\texttt{["M2"]["inject\_time"]}}, and {\color{green}\texttt{["M3"]["inject\_time"]}}.
  \item {\color{cyan}$D_1=0.01\;\textrm{mm}^2\cdot\textrm{h}^{-1}$}; accessed by {\color{green}\texttt{["Vi"]["diffusion"]}}.
  \item {\color{cyan}$\gamma=0.013\;\textrm{h}^{-1}$}; accessed by {\color{green}\texttt{["Vi"]["grow\_neg"]}}.
  \item {\color{cyan}$b=300\;\textrm{virus}\cdot\textrm{cell}^{-1}$}; accessed by {\color{green}\texttt{["Vi"]["burst\_h"]}}.
\end{itemize}

{\bf Boundary conditions}. 
Due to the radial symmetry assumption, we adopt the no-flux conditions for cell number densities and virus at $r=0$:
\begin{align}
  \label{eq:eqn_pde_cbc}
  &\left.\pp{B}{r}\right|_{r=0} = 
  \left.\pp{H}{r}\right|_{r=0} = 
  \left.\pp{N}{r}\right|_{r=0} = 
  \left.\pp{Q}{r}\right|_{r=0} = 
  \left.\pp{M_k}{r}\right|_{r=0} = 
  \left.\pp{M_k^0}{r}\right|_{r=0} = 0\;,\quad k=1,2,3\;, \\
  \label{eq:eqn_pde_vbc}
  &\left.\pp{V}{r}\right|_{r=0} = 0\;;
\end{align}
to this end, the advective velocities are naturally zero at this location: $U_1(0,t)=U_2(0,t)=U_3(0,t)=0$.
The mean velocity $E(r,t)$ is computed by integrating~(\ref{eq:eqn_pde_vel_e}) from the tumor center to $r$, with the boundary condition:
\begin{equation}\label{eq:eqn_pde_ebc}
  E(0,t) = 0\;.
\end{equation}
Because the tumor boundary is defined by the mean velocity, no boundary condition is required for cells that are entirely transported by $E$, i.e., $B$, $H$, $Q$, $N$, and $M_k^0$, $k=1,2,3$.
The immune cells $M_k$, $k=1,2,3$ enters the tumor through the boundary, which leads to a Dirichlet type condition at $r=R(t)$:
\begin{equation}\label{eq:eqn_pde_rbc}
  M_1(R(t),t) = M_2(R(t),t) = M_3(R(t),t) = \frac{1}{3}M_{\textrm{bc}}(t)\;,
\end{equation}
where $M_{\textrm{bc}}(t)$ is a prescribed function that is potentially time-dependent and it designates the rate at which immune cells enter the tumor.
For now, we consider a simple constant function defined as follows:
\begin{equation}\label{eq:eqn_pde_tcell_bcval}
  {\color{cyan}M_{\textrm{bc}}(t) = 0.05}\;,
\end{equation}
that is, the tumor is exposed to immune cells with concentration $5\%$ at all times (they will not enter the tumor, however, unless there is a gradient of the corresponding driving cell type at tumor boundary).

It is assumed that virus stays inside the tumor, hence homogeneous Dirichlet boundary condition for virus concentration is applied at the tumor boundary:
\begin{equation}\label{eq:eqn_pde_virus_dbc}
  V(R(t),t) = 0\;.
\end{equation}

{\bf Initial conditions}.
Lastly, we assume at the beginning, all cells are melanoma tumor cells, hence for all $0\le r\le R_0$:
\begin{equation}\label{eq:eqn_pde_cell_ic}
  B(r,0) = 1\;,\ 
  H(r,0) = Q(r,0) = N(r,0) = M_k(r,0) = M_k^0(r,0) = 0\;,\ 
  k = 1,2,3,
\end{equation}
where {\color{cyan}$R(0) = R_0 = 4\textrm{mm}$} is the initial radius of the tumor.

The virus is injected at the tumor center at $t=0$, which gives rise to the initial condition:
\begin{equation}\label{eq:eqn_pde_virus_ic}
  V(r,0) = \frac{V_0}{4\pi r^2}\delta_0(r)\;,\quad0\le r\le R_0\;.
\end{equation}
for the virus concentration.
Here {\color{cyan}$V_0 = 10^4\textrm{virus}$} is the total amount of injected viruses.

\section{A model problem testing point source}
\label{sec:model}
{\bf The governing equations}.
It is possibly a simplest model that involves a moving domain, a point source, and an elastic interstitial space. 
\begin{subequations}\label{eq:model_pde}
  \begin{align}
    \label{eq:model_pde_m}
    &\pp{M}{t} + \frac{1}{r^2}\po{r}\left[r^2MV\right] = \lambda M + h(r,t)\;,\\
    \label{eq:model_pde_s}
    &\pp{S}{t} + \frac{1}{r^2}\po{r}\left[r^2SV\right] = -\kappa(S-S_0)-h(r,t)\;,\\
    \label{eq:model_pde_v}
    &\frac{1}{r^2}\po{r}\left[r^2 V\right] = \lambda M-\kappa(S-S_0)\;.
  \end{align}
\end{subequations}
For simplicty, we assume the initial condition is given by $M\equiv1-S_0$ and $S\equiv S_0$ on the domain $r\in[0,\;R_0]$, $R_0=10$.
No-flux condition is enforced for both $M$ and $S$ at $r=0$, and no boundary condition is needed at $r=R(t)$.
Note that if $h(r,t)\equiv0$, we have $R(t) = R_0\exp(\lambda(1-S_0)t/3)$.

To test the source term and the elastic term, we shall consider:
\begin{equation}\label{eq:model_source}
  h(r,t) = n_M\delta(r-r_0)\delta(t-t_0)\;,
\end{equation}
where $r_0=6$ and $t_0=10$.

\end{document}

